\section{Proof details}
\label{app:proofs}

\subsection{Boolean interval recurrence}

For \(S\subseteq[r]\), the interval \([\varnothing,S]\) in \(2^{[r]}\) is
again Boolean.  The incidence recurrence
\begin{equation}
 \sum_{T\subseteq S}\mu(\varnothing,T)=\delta_{S,\varnothing}
 \label{eq:boolean-recurrence}
\end{equation}
is solved by \(\mu(\varnothing,T)=(-1)^{|T|}\), since
\(
 \sum_{k=0}^{|S|}\binom{|S|}{k}(-1)^k=(1-1)^{|S|}
\).
This proves the sign used in \eqref{eq:boolean-weight} without importing a
numerical label or factorization oracle.

If \(a_1,\ldots,a_r\) are distinct prime covers in \(\MZ\), their full join
is \(\prod_i a_i\).  Every divisor of that squarefree product is a unique
subset product, so \(\beta_A\) is a pointed order isomorphism.  If the atoms
are coordinate unit vectors in \(\Ffree\), the same proof replaces subset
products by \(0/1\) exponent vectors.

\subsection{Cumulant coefficient identity}

Grouping the partitions in \cref{lem:factorization-cancellation} by their
number of blocks gives
\begin{equation}
 c_r=\sum_{k=1}^r S(r,k)(-1)^{k-1}(k-1)!,
 \label{eq:stirling-cumulant}
\end{equation}
where \(S(r,k)\) is a Stirling number of the second kind.  Using
\(
 \sum_{r\ge k}S(r,k)z^r/r!=(e^z-1)^k/k!
\), one obtains
\begin{align}
 \sum_{r\ge1}c_r\frac{z^r}{r!}
 &=\sum_{k\ge1}\frac{(-1)^{k-1}}{k}(e^z-1)^k\\
 &=\log(1+e^z-1)=z.
 \label{eq:cumulant-egf}
\end{align}
Thus \(c_1=1\) and all higher coefficients vanish.  The argument applies at
every arity, although the executable candidate froze only ranks two and
three.

\subsection{Intervals and M{\"o}bius values under valuation}

For \(m\mid n\), valuation restricts to
\begin{equation}
 [m,n]_{\mid}\ \iso\
 [\PhiVal(m),\PhiVal(n)]_{\le}.
 \label{eq:interval-isomorphism}
\end{equation}
Indeed, \(m\mid d\mid n\) is equivalent coordinatewise to
\(
 \PhiVal(m)\le\PhiVal(d)\le\PhiVal(n)
\), and unique factorization supplies the inverse.  Applying the incidence
recurrence on the two isomorphic intervals gives
\begin{equation}
 \mu_{\MZ}(m,n)
 =\mu_{\Ffree}(\PhiVal(m),\PhiVal(n)).
 \label{eq:mobius-transport}
\end{equation}
All constructions made from interval convolution therefore transport.

\subsection{Decoration transport and global invariants}

Let \(T_X\) be any source-resident \(k\)-index tensor.  Define its clone by
\begin{equation}
 T_{\Ffree}(\PhiVal(x_1),\ldots,\PhiVal(x_k))
 :=T_{\MZ}(x_1,\ldots,x_k).
 \label{eq:tensor-transport}
\end{equation}
The same rule transports roof marks, Gram kernels, compiled coefficient
arrays, and compatible active cutoffs.  A global invariant may aggregate
over all entries and all cutoff levels; equation~\eqref{eq:tensor-transport}
still supplies an isomorphism of the complete filtered decorated object.
The proof of \cref{thm:clone} therefore does not reduce a nonlocal
construction to a local one.

\subsection{Uniform majorant on compact substrips}

Let \(K\) be a compact subset of
\(
 1-2\eta<\Ree s<2\eta
\).  There is \(\varepsilon>0\) such that for all \(s\in K\),
\begin{equation}
 2\eta+\Ree s\ge1+\varepsilon,
 \qquad
 2\eta+1-\Ree s\ge1+\varepsilon.
 \label{eq:uniform-exponents}
\end{equation}
Consequently the absolute value of \eqref{eq:mixed-functional} is bounded,
uniformly on \(K\), by a constant multiple of
\begin{equation}
 \left(\sum_{n\ge2}n^{-1-\varepsilon}\right)^2.
 \label{eq:uniform-majorant}
\end{equation}
This proves normal convergence independently of prime-density estimates.
The result is a statement about the new scalar functional only; it supplies
no Schatten or trace-class classification of an unconstructed operator.

\subsection{Why exponentiation does not upgrade the object}

Given any holomorphic \(h(s)\), the function \(\exp(h(s))\) is holomorphic
and zero-free.  This elementary fact does not make it a Fredholm determinant.
A determinant requires an operator family and the corresponding trace-ideal
or relative hypotheses.  Similarly, multiplying the inherited
\(\detthree(I-z\cB_s)\) by \(\exp[-z^2\MK(s)/2]\) would declare a new
scheme-dependent product.  It would not restore the quadratic term to the
logarithm of \(\detthree\) or convert \(\MK\) into an ordinary trace.
